% RESEARCH_DRAFT_SHA256: 27e676e1aaa36fe63c9fe446f3718414ec2630b8a49e803a69c8df8f8641b760
% Layout and source-reference wrappers expanded for the web renderer.
\documentclass[11pt]{article}
\usepackage[margin=25mm]{geometry}
\usepackage{amsthm}
\usepackage{fontspec,unicode-math}
\setmainfont{Latin Modern Roman}
\setmathfont{Latin Modern Math}
\usepackage{amsmath,amssymb,mathtools}
\usepackage{xurl,hyperref}
\hypersetup{colorlinks=true,urlcolor=blue,linkcolor=blue}
\setcounter{secnumdepth}{0}
\setlength{\parindent}{0pt}
\setlength{\parskip}{5pt}
\setlength{\emergencystretch}{3em}
\newcommand{\sref}[2]{\hyperlink{source:#1:#2}{[S#2]}}
\newcommand{\cataloguescope}[1]{\paragraph{Source question.}#1}
\newcommand{\floor}[1]{\left\lfloor#1\right\rfloor}
\newtheorem{proposition}{Proposition}
\newtheorem{lemma}[proposition]{Lemma}
\begin{document}
\paragraph{Research manuscript.} OpenAI \texttt{gpt-6-astra}, reasoning \texttt{ultra}; 11 October 2026.
\paragraph{Preserved source body.} The complete catalogue mathematical text and sources follow with layout and source links expanded, including source qualifications and its historical preparation notes. New research begins at the section headed Research attempt.

\section{Minimal Convex Difference Decompositions of CPWL Functions}\label{statement}

\subsection{Definitions and mathematical statement}

\paragraph{Shared notation.}Standard number systems are denoted by $\mathbb N,\mathbb Z,\mathbb Q,\mathbb R,\mathbb C$; index ranges, computational models and specialized notation follow the cited problem source and the definitions below.


For continuous piecewise-linear $f$, each affine region carries a linear formula. A convex-difference representation is $f=g-h$ with convex piecewise-linear $g,h$; minimizing pieces requires counting the regions according to the source\textquotesingle s convention.

A continuous piecewise-linear function is continuous and agrees with an affine function on each member of a finite polyhedral partition. A difference-of-convex decomposition writes it as $g-h$ with both $g$ and $h$ convex.
\paragraph{Statement recorded in the supplied source.}
Given a continuous piecewise-linear function $f$, decompose it as $f=g-h$ with $g$ and $h$ convex continuous piecewise-linear functions having as few linear pieces as possible. Construct a finite procedure that finds a globally minimal decomposition in arbitrary dimension. \href{https://doi.org/10.4171/owr/2025/12}{[S2]}
\paragraph{Scope qualification.}
The objective counting $g$ and $h$ pieces, encoding, and treatment of compatible subdivisions are not specified. The target is global across decompositions, not optimality on one fixed complex.
\paragraph{Source status.}
The supplied catalogue labels this record partially solved. Its source-status metadata was checked against the supplied snapshot on 2026-10-11. This import does not independently establish current open status or mathematical correctness. The catalogue literature assessment is dated 2026-08-22.
\subsection{Short English statement}
Find a finite globally optimal procedure for decomposing a piecewise-linear function as a difference of two convex ones.
\subsection{Sources}
\begin{itemize}
\item[S1] Ulam AI and the UnsolvedMath Contributors, UnsolvedMath: A Curated Collection of Open Mathematics Problems, catalogue entry. Catalogue curation and metadata: CC BY 4.0. The supplied source question is retained; definitions and the English summary are editorial additions.
\url{https://huggingface.co/datasets/ulamai/UnsolvedMath}.
\item[S2] Mini-Workshop: Alcoved Polytopes in Physics and Optimization (2025), Oberwolfach Reports, DOI 10.4171/owr/2025/12, pp. 531--558
\url{https://doi.org/10.4171/owr/2025/12}.
\item[S3] Decomposition polyhedra of CPWL functions
\url{https://openreview.net/forum?id=vVCHWVBsLH}.
\textit{Underlying Oberwolfach Reports material retains its source-specific licence, including CC BY-SA 4.0 where applicable; see THIRD\_PARTY\_NOTICES.md.}
\end{itemize}
\paragraph{Preparation and review.}
Original question attributed to the cited source and UnsolvedMath contributors. Codex (GPT-6; exact serving revision not exposed) assisted with definitions, English summaries, source inspection and formatting. No new human mathematical authorship or independent proof verification is asserted. Catalogue identity checked against the complete supplied registry and existing repository statements; known flagged repetitions and obvious identical conjectures excluded. Full mathematical-equivalence/literature review has not been performed. Statement preparation and final review followed contribution prompts C and D; automated checks verify structure and identity, not open status.
\label{end:statement}

\section{Research attempt}
\subsection{Literature, source context and scope}
The full original catalogue body and source qualifications are preserved above. The relevant primary workshop passages were read in the cached publisher report. Historical claims and subsequent attributed results are distinguished from the mathematical work below. A restricted example, standard lemma or known theorem is not counted as new progress on the entire original target. Missing assumptions and repeated or rhetorical extracts are not silently turned into different solved problems.
Codex (GPT-6; exact serving revision not exposed) supplied the auxiliary proof route used here. OpenAI \texttt{gpt-6-astra}, reasoning \texttt{ultra}, independently worked through its mathematical steps, wrote this exposition, and performed the finite checks. This is model collaboration, not independent human verification.

The workshop asks for globally minimal convex-difference decompositions, beyond those compatible with a fixed complex. The September 2026 publisher version retains that distinction in Problem 1.1 and Section 1.1. This attempt gives a complete one-dimensional auxiliary calculation. It is a classical slope decomposition, not new progress in general dimension. The root supplied the same auxiliary route; it was independently checked here.
\subsection{Global piece minima on the real line}
Let $f:\mathbb R\to\mathbb R$ be continuous piecewise affine with finitely many actual breakpoints $t_1<\cdots<t_m$, nonzero slope jumps $\Delta_j$, and affine formula $ax+b$ to the left of all breakpoints. Let $p$ count positive jumps and $q$ negative jumps, so $p+q=m$. Count pieces as maximal intervals on which a convex CPWL function is affine; an affine function has one piece.
\begin{proposition}
Among all representations $f=g-h$ by convex CPWL functions, the coordinatewise minimum piece counts are $(p+1,q+1)$, attainable simultaneously. In particular the minimum sum of piece counts is $m+2$.
\end{proposition}
\begin{proof}
Matching the initial affine formula and successive slopes gives
\[
 f(x)=ax+b+\sum_{j=1}^m\Delta_j(x-t_j)_+.
\]
Indeed the right side has the same slope on every intervening interval and agrees on the first interval; continuity then gives equality everywhere. Set
\[
 g(x)=ax+b+\sum_{\Delta_j>0}\Delta_j(x-t_j)_+,
 \qquad h(x)=\sum_{\Delta_j<0}(-\Delta_j)(x-t_j)_+.
\]
Each hinge has nondecreasing slope, so both functions are convex, and their strictly positive jumps occur at exactly $p$ and $q$ distinct positions. Thus they have $p+1$ and $q+1$ pieces.

For any other convex CPWL pair $G,H$, the one-sided derivatives exist and their jumps $\alpha_j,\beta_j$ at $t_j$ are nonnegative. Since $G-H=f$, $\alpha_j-\beta_j=\Delta_j$. A positive $\Delta_j$ forces a positive jump of $G$, and a negative one forces a positive jump of $H$. Hence $G$ has at least $p$ distinct breakpoints and $H$ at least $q$. Additional breakpoints outside this list cannot reduce either count. This proves the unrestricted one-dimensional lower bounds.
\end{proof}
For example $f(x)=x_+-2(x-1)_++(x-2)_+$ has minimal counts $(3,2)$. For rational sorted breakpoint/slope data the construction uses $O(m)$ arithmetic operations; sorting unsorted distinct positions costs $O(m\log m)$ comparisons. This is an arithmetic count, not a claim that bit lengths are constant.
\section{Verification and remaining gaps}
Exact checks compare the hinge formulas interval by interval and count the jump signs, including the affine case $m=0$. In higher dimensions, independently repairing jumps across faces need not yield globally compatible convex functions. This proof supplies neither an enumeration bound for unknown complexes nor a global minimum algorithm in arbitrary dimension.

\section{References and source attribution}
Primary workshop report: \url{https://doi.org/10.4171/owr/2025/12}. Relevant passages read: OWR 2025/12, printed pages 551--553, Brandenburg Problem1 and fixed-complex theorem.

\url{https://link.springer.com/article/10.1007/s00454-026-00875-1}. Read published September2026 abstract, Problem1.1 and Sections1.1.1/1.1.6: fixed-complex versus global distinction remains explicit; full paper not audited.

All retrievals were checked on 11 October 2026. Original question attribution remains in the preserved Sources section. External theorems are attributed and not claimed as new discoveries or fully verified merely by citation.


\paragraph{Integration review.} Codex (GPT-6; exact serving revision not exposed) reviewed the written arguments and expanded layout and source-reference wrappers. Repository metadata and the production web renderer were checked using Python 3.13 and JavaScript with MathJax 3.2.2. This model-assisted review is not independent human proof verification; no PDF compilation or formal proof checking is asserted.

\section{Completion Estimate}
\noindent\textbf{COMPLETION: 0\%}
No substantiated new progress toward the entire original target is claimed. The worked results and failed route are delimited above; a known result, illustrative model, or auxiliary obstruction is not counted as original-target completion.

\section{Final status and attribution}
\textbf{Attempt status: unresolved.} This describes this research attempt, not a new ruling on the source catalogue's status or on external literature claims.
The mathematics and this self-audit were generated by OpenAI \texttt{gpt-6-astra}, reasoning \texttt{ultra}, on 11 October 2026. Codex (GPT-6; exact serving revision not exposed) orchestrated the batch, read and reviewed the written mathematical arguments, independently reran the finite checks, and contributed formatting/repository review. No human mathematical authorship, independent human proof checking, or formal verification is claimed. The historical preparation paragraphs in the reproduced catalogue body concern that body only.

Contribution prompt A was applied to the complete definition and supplied source record, followed by the adversarial checks of prompt B. Local reproduction files are \texttt{batch026/checks.py}, \texttt{computations.json}, and the per-ID audit JSON. They were not public links at preparation time. Numerical checks are finite and do not prove unrestricted results. No LaTeX engine was found locally; compilation was not performed by this research agent.

\end{document}
